\documentclass[main.tex]{subfiles}


\begin{document}

%from pagenumber to up
% \phantomsection 
% \hypertarget{toc}{}

 

 

% номер секции вручную
\setcounter{section}{37
}
\addtocounter{section}{-1} 

\section{Механическая волна}


\textbf{Механическая волна}~--- это распространение колебаний частиц в упругой среде (твёрдой, жидкой или газообразной). Выделяют два вида волн.
\begin{enumerate}
    \item \textbf{Поперечная}: частицы среды колеблются \emph{перпендикулярно} направлению распространения волны (распространяется только в \emph{твердых средах}). Такая волна показана на рис.\,\ref{fig:p62}; направление колебаний слоев (вертикальных рядов частиц) перпендикулярно \emph{скорости волны}~$\vec v$.

    \begin{figure}[H]
    \centering

    \begin{tikzpicture}[font=\small,            line cap=round,                             >={Stealth[inset=0pt,round,length=3mm, angle'=20]},vec/.style={>={Stealth[inset=0pt,round,length=3.5mm, angle'=20]}}]

%lambda
\draw[densely dashed,gray] ({360/4.5 pt},{.3*cos(360)}) --++ (0,.4) coordinate (l1);
\draw[densely dashed,gray] ({720/4.5 pt},{.3*cos(360)}) --++ (0,.4) coordinate (l2);

\draw[<->,gray] ([yshift=-.2cm]l1) --++ ({360/4.5 pt},0) node[midway,below,black] {$\lambda$};

%v
\draw[->,vec,thick,NavyBlue] ({1440/4.5 pt},{(-36/4.5*4pt)-.3cm}) ++ (0,-0.4) ++ (-1,0) --++ (1,0) node[midway,above,black] {$\vec v$};

\foreach \y in {0,1,2,3,4}{
    \foreach \x in {0,36,...,1440}
    {
    \shade[ball color=orange,yshift={-36/4.5*\y pt}] ({\x/4.5 pt},{.3*cos(\x)}) circle (.05);
    }
}

    \end{tikzpicture}
\caption{Поперечная волна}
\label{fig:p62}
\end{figure}
\nointerlineskip
\item \textbf{Продольная}: частицы среды колеблются \emph{параллельно} направлению распространения волны. Такая волна показана на рис.\,\ref{fig:p63}; направление колебаний слоев (вертикальных рядов частиц) параллельно скорости волны~$\vec v$.

    \begin{figure}[H]
    \centering

    \begin{tikzpicture}[font=\small,            line cap=round,                             >={Stealth[inset=0pt,round,length=3mm, angle'=20]},vec/.style={>={Stealth[inset=0pt,round,length=3.5mm, angle'=20]}}]

%lambda
\draw[densely dashed,gray] ({360/4.5 pt},0) --++ (0,.5) coordinate (l1);
\draw[densely dashed,gray] ({720/4.5 pt},0) --++ (0,.5) coordinate (l2);

\draw[<->,gray] ([yshift=-.2cm]l1) --++ ({360/4.5 pt},0) node[midway,above,black] {$\lambda$};

%v
\draw[->,vec,thick,NavyBlue] ({1440/4.5 pt},{(-36/4.5*4pt)-.3cm}) ++ (0,-0.4) ++ (-1,0) --++ (1,0) node[midway,above,black] {$\vec v$};


\foreach \z in {0,1,2,3}{
\foreach \y in {0,1,2,3,4}{
    \foreach \x in {0,36,54,...,144}
    {
    \begin{scope}[xshift={2*180pt/4.5*\z}]
    \shade[ball color=orange,yshift={-36/4.5*\y pt}] ({(180pt/4.5-180pt/4.5*cos(\x))},0) circle (.05);
    \end{scope}
    }
}
}

\foreach \z in {0,1,2,3}{
\foreach \y in {0,1,2,3,4}{
    \foreach \x in {23,157}
    {
    \begin{scope}[xshift={2*180pt/4.5*\z}]
    \shade[ball color=orange,yshift={-36/4.5*\y pt}] ({(180pt/4.5-180pt/4.5*cos(\x))},0) circle (.05);
    \end{scope}
    }
}
}

%last seris
\foreach \y in {0,1,2,3,4}{
    \foreach \x in {180}
    {
    \begin{scope}[xshift={2*180pt/4.5*3}]
    \shade[ball color=orange,yshift={-36/4.5*\y pt}] ({(180pt/4.5-180pt/4.5*cos(\x))},0) circle (.05);
    \end{scope}
    }
}


    \end{tikzpicture}
\caption{Продольная волна}
\label{fig:p63}
\end{figure}
\end{enumerate}
\nointerlineskip

\textbf{Длина волны} ($\lambda$ [м])~--- это расстояние, на которое распространяется волна за период колебаний:
\begin{equation}
    \lambda=vT.
\end{equation}
Длина волны~$\lambda$ в поперечной волне равна расстоянию между \emph{соседними} горбами или впадинами (рис.\,\ref{fig:p62}). В продольной волне~$\lambda$ есть расстояние также между \emph{соседними} сжатиями или разрежениями (рис.\,\ref{fig:p63}). 
% (Вообще говоря, длиной волны называется расстояние (вдоль направления распространения волны, указываемого скоростью волны~$\vec v$) между двумя ближайшими частицами среды, колеблющимися в одинаковой фазе (то есть частицы имеют одинаковые положения, скорости и ускорения вдоль направления их колебаний.)

\textbf{Звук}~--- это механическая волна в диапазоне частот от 20~Гц до 20~кГц. Ниже этого диапазона лежит область \emph{инфразвука}, выше~--- область \emph{ультразвука}.

\sbox{0}{%
    \begin{tikzpicture}[font=\small,            line cap=round,                             >={Stealth[inset=0pt,round,length=3mm, angle'=20]},vec/.style={>={Stealth[inset=0pt,round,length=3.5mm, angle'=20]}}]


\filldraw[lightgray,semitransparent] (5,2.5) rectangle (5.3,0.5);
\draw (5,2.5) -- (5,0.5);

% А
\draw [dashed,orange] (3,2) sin (4,2.1) node[above,black] {А} cos (5,2);
\draw [dashed,orange] (3,2) -- (5,2);
\draw [orange,thick] (3,2) sin (4,1.9) cos (5,2);

\node at (3,2) [circle,fill=black,inner sep=0pt,minimum size=1mm,label=below:] {};
\node at (5,2) [circle,fill=black,inner sep=0pt,minimum size=1mm,label=below:] {};

% Б
\begin{scope}[yshift=-1cm]
\draw [thick,orange] (4,2) sin (4.5,2.2) cos (5,2);
\draw [dashed,orange] (4,2) -- (5,2);
\draw [orange,dashed] (4,2) sin (4.5,1.8) node[below,black] {Б} cos (5,2);

\node at (4,2) [circle,fill=black,inner sep=0pt,minimum size=1mm,label=below:] {};
\node at (5,2) [circle,fill=black,inner sep=0pt,minimum size=1mm,label=below:] {};
\end{scope}



    \end{tikzpicture}%
}

{%
\setlength\intextsep{0pt}
\begin{wrapfigure}[17]{r}{\wd0}
    \centering
    \usebox{0}%
    \caption{Струны}
    \label{fig:p64}
\end{wrapfigure}

Вот основные характеристики звука.
\begin{itemize}
    \item \textbf{Громкость} звука определяется \emph{амплитудой} колебаний. Так, при колебаниях струн (рис.\,\ref{fig:p64}) громкость струны~Б  больше, чем громкость струны~А (что видно по <<размаху>>).
    \item \textbf{Высота} звука определяется \emph{частотой} колебаний. Например, в опыте на рис.\,\ref{fig:p64} при различии только в длинах струн звук струны~Б кажется выше, чем звук струны~А.
    % звук струны~Б кажется выше, чем звук струны~Б (при различии только в их длинах, что обнаружено на опыте).
\end{itemize}

}

\emph{Скорость звука} в твердых телах больше, чем в жидкостях, а в жидкостях~--- больше, чем в газах. Следует отметить, что \emph{частота} звука при переходе из одной среды в другую не меняется и определяется частотой источника звука. Сказанное относится к любым механическим волнам. 



























 
\end{document}